% ==============================================================
\section{Conclusion}\label{sec:conclusion}
% ==============================================================

This paper formally defined the \textbf{event attribution problem}
and constructed a five-step attribution pipeline as the first statistical solution.

Experiments yielded the following findings:
\begin{itemize}
\item High attribution accuracy was achieved across diverse data conditions including
Zipf distributions; under null conditions over 100 seeds,
a mean FDR of 0.070 ($\leq \alpha = 0.10$) was confirmed, demonstrating FDR control (EX1).
\item Ablation analysis demonstrated the complementarity of proximity and magnitude,
and the progressive false positive reduction effect of global BH correction
and deduplication (EX2).
\item Comparison against 4 related methods confirmed that the proposed method
achieves the best or equal F1 across all 5 conditions (mean F1 = 0.68),
with particularly prominent advantage under confounding and overlap conditions (EX3).
\item Exploratory application to real-world data (Dunnhumby, 276K baskets)
confirmed attribution generation for all campaigns (100\% coverage),
demonstrating applicability (EX4).
\end{itemize}

\subsection{Limitations and Future Work}

This pipeline has the following limitations.
First, accuracy degrades under conditions of overlapping or high-density event windows.
Second, the circular shift permutation test assumes weak stationarity;
data with strong trends or seasonality may require extensions such as block permutation.
Third, there is a tradeoff between detection power and weak signal recovery.
Future work includes handling non-stationary time series
and evaluating external validity on large-scale real-world data.
